% GENERATED FILE. Edit canonical lessons, then run npm run generate:books.
% canonical-source-hash: fnv1a64:d089d3b2347cec6d
% canonical-lessons: MR-R18-script-warmup, MR-R18-script-a-r4, MR-R18-script-b-r4, MR-R18-script-c-r4, MR-R18-script-d-r4, MR-R22-marks-r4, MR-R22-two-rows-r4, MR-R22-curled-r4, MR-R22-sibilants-r4

\chapter{The Doorway Signs at Long Distance}
\label{ch:doorway-script-distance}
\hlchaptermodality{22}

\begin{chapteropening}
\textbf{By the end of this chapter:} \emph{I can write dhanyavad and retrieve all twelve doorway signs after more than eighty later lessons.}

Retrieve the final sign batch from sound after every new doorway sign has crossed the long-distance window.
\end{chapteropening}

\section[dhanyavād]{The doorway word, much later}

\label{lesson:MR-R18-script-warmup}



\begin{quote}
Close every earlier model and listen before picking up the pencil.
\end{quote}

\begin{tcolorbox}[breakable,colback=teal!4,colframe=teal!35!black,title={Before you move on}]
Hear \emph{dhanyavād}. Say its meaning, read \textbf{\mr{धन्यवाद}}, then cover the model and write it once. This reconnects meaning and script without adding anything new.
\end{tcolorbox}
\section[visarga / aa / bha]{First batch at long distance}

\label{lesson:MR-R18-script-a-r4}



\begin{quote}
Keep every model closed and one blank line ready.
\end{quote}

\begin{tcolorbox}[breakable,colback=teal!4,colframe=teal!35!black,title={Before you move on}]
From sound and spoken names only, write \emph{dhanyavād}, visarga, independent \emph{ā}, and \emph{bha}. Compare once: \textbf{\mr{धन्यवाद}; \mr{ः} \mr{आ} \mr{भ}}.
\end{tcolorbox}
\section[e / anusvara / ta]{Second batch at long distance}

\label{lesson:MR-R18-script-b-r4}



\begin{quote}
Keep every model closed and one blank line ready.
\end{quote}

\begin{tcolorbox}[breakable,colback=teal!4,colframe=teal!35!black,title={Before you move on}]
From spoken names only, write e-mark, anusvāra, and dental \emph{ta}. Compare once: \textbf{\mr{े} \mr{ं} \mr{त}}.
\end{tcolorbox}
\section[da / dha / ba]{Voiced batch at long distance}

\label{lesson:MR-R18-script-c-r4}



\begin{quote}
Keep every model closed and one blank line ready.
\end{quote}

\begin{tcolorbox}[breakable,colback=teal!4,colframe=teal!35!black,title={Before you move on}]
From spoken names only, write \emph{da}, \emph{dha}, and \emph{ba}. Compare once: \textbf{\mr{द} \mr{ध} \mr{ब}}.
\end{tcolorbox}
\section[ya / lla / va]{Final batch at long distance}

\label{lesson:MR-R18-script-d-r4}



\begin{quote}
Keep every model closed and one blank line ready.
\end{quote}

\begin{tcolorbox}[breakable,colback=teal!4,colframe=teal!35!black,title={Before you move on}]
From spoken names only, write \emph{ya}, curled-back \emph{ḷa}, and \emph{va}. Compare once: \textbf{\mr{य} \mr{ळ} \mr{व}}.
\end{tcolorbox}
\section[sahā khuṇā]{The six marks, a whole book later}

\label{lesson:MR-R22-marks-r4}



\begin{quote}
You have not been asked for these since the counting chapter. Say them anyway, and notice which one comes slowest.
\end{quote}

\begin{tcolorbox}[breakable,colback=teal!4,colframe=teal!35!black,title={Before you move on}]
Short \emph{i} — before, above, or below? (\textbf{Before.}) {[}PAUSE 2s{]} Short \emph{u}? (\textbf{Below.}) {[}PAUSE 2s{]} Long \emph{ū}? (\textbf{Below, deeper.}) {[}PAUSE 2s{]} The vocalic \emph{r}? (\textbf{Below, opening the other way.}) {[}PAUSE 2s{]} The candrabindu? (\textbf{Above, in a crescent.}) {[}PAUSE 2s{]} And short \emph{a} standing alone? (\textbf{Its own letter.})
\end{tcolorbox}
\section[don olī]{Two rows, at book distance}

\label{lesson:MR-R22-two-rows-r4}



\begin{quote}
If the pattern held, you will not be recalling seven separate things. Find out.
\end{quote}

\begin{tcolorbox}[breakable,colback=teal!4,colframe=teal!35!black,title={Before you move on}]
Velar row, plain to voiced-and-breathy. (\emph{\textbf{ka, kha, ga, gha\emph{.\textbf{) Palatal row, the same. (}}ca, cha, ja, jha\emph{.\textbf{) {[}PAUSE 2s{]} Which sign does \emph{mājhe} need? (}}jha\emph{.\textbf{) {[}PAUSE 2s{]} Which one opens \emph{gheṇe}? (}}gha\emph{.\textbf{)}}}}
\end{tcolorbox}
\section[vaḷlelī olī]{The curled row, a whole book later}

\label{lesson:MR-R22-curled-r4}



\begin{quote}
These five are the reason the romanizations in this book carry dots under some letters and not others.
\end{quote}

\begin{tcolorbox}[breakable,colback=teal!4,colframe=teal!35!black,title={Before you move on}]
The four curled letters. (\emph{\textbf{ṭa, ṭha, ḍa, ṇa\emph{.\textbf{) {[}PAUSE 2s{]} What is the dot under \emph{ṭ} telling you? (}Curl the tongue back.\textbf{) {[}PAUSE 2s{]} Which letter ends \emph{bolṇe} and \emph{khāṇe}? (}}ṇa\emph{.\textbf{) {[}PAUSE 2s{]} And \emph{ba} with the voice switched off? (}}pa\emph{.\textbf{)}}}}
\end{tcolorbox}
\section[śevaṭce sahā]{The last six, at book distance}

\label{lesson:MR-R22-sibilants-r4}



\begin{quote}
The end of the second runway, asked for once more with everything else in between.
\end{quote}

\begin{tcolorbox}[breakable,colback=teal!4,colframe=teal!35!black,title={Before you move on}]
The plain \emph{l}, beside the curled one. (\emph{\textbf{la\emph{.\textbf{) {[}PAUSE 2s{]} The two sibilants most speakers say alike. (}}śa\emph{ and }ṣa\emph{.\textbf{) {[}PAUSE 2s{]} The three vowels that stand alone as well as leaning. (}}u\emph{, }ū\emph{, }e\emph{.\textbf{) {[}PAUSE 2s{]} Why is a letter still written for a sound that merged? (}Spelling outlives the sound it recorded.\textbf{)}}}}
\end{tcolorbox}
